\chapter{From Symbolic to Scalar: The Main Binding Theorem}
\label{ch:specialization}

The symbolic protocol checks polynomial identities in $T$. Evaluating every word at a
single point $T=c$ turns a symbolic transcript into a scalar one with challenges
$\beta_i=cr_i$ (\cref{lem:specialization}). This chapter makes the reduction exact: an
analysis-only point $c$, uniform and independent of the protocol, couples the symbolic
game with a scalar game, so that the master path lemma applies at $T=c$; a root count
in $c$ then lifts the scalar conclusion to an identity of polynomials in $T$.

%=====================================================================
\section{The symbolic game}
%=====================================================================

\begin{definition}[Symbolic game]\label{def:symbolic-game}
A \emph{symbolic prover} $\mathsf P^\ast$ is an arbitrary (unbounded, possibly
randomized) strategy for \cref{prot:sfri}, with the following conventions.
\begin{enumerate}[label=\textup{(\roman*)}]
\item $W_0:D_0\to\F$ (leaves of $T$-degree $0$; non-constant leaves are rejected).
\item $W_{i}:D_{i}\to\K[T]$ for $1\le i\le k$, with \emph{no} restriction on the $T$-degree.
\item $\gt\in\K[T]$ with $\deg\gt\le d$ (otherwise $\mathsf V$ rejects).
\end{enumerate}
The \emph{symbolic commit transcript} is
$\tau:=(W_0,r_1,W_1,\dots,r_k,W_k,\gt)$; $r:=(r_1,\dots,r_k)$. The symbolic passing sets
$P_i^{\mathrm{sym}},Q_{i+1}^{\mathrm{sym}}$ and $\pi^{\mathrm{sym}}_i$ are defined as in
\cref{def:passing}, with every equality read in $\K[T]$, $\Fold_{\beta_{i+1}}$ replaced by
$\Fold_{r_{i+1}T}$, and $g$ replaced by $\gt$. $\Acc$ denotes acceptance of
\cref{prot:sfri}.
\end{definition}

Allowing arbitrary $T$-degrees in (ii) only strengthens the prover; the analysis below
never uses a degree bound on intermediate layers.

\begin{lemma}[Passing sets in the symbolic game]\label{lem:passing-sym}
\cref{lem:passing} holds verbatim for the symbolic game; in particular
$\Pr[\Acc\mid\tau]=(\pi^{\mathrm{sym}}_0)^s$ and $\pi^{\mathrm{sym}}_i\le\pi^{\mathrm{sym}}_{i+1}$.
\end{lemma}

\begin{proof}
The proof of \cref{lem:passing} uses only the shape of the checks, not the value ring.
\end{proof}

\begin{definition}[Specialized transcript]\label{def:specialized}
For $c\in\K$ and a symbolic commit transcript $\tau$, the \emph{specialization at $c$} is
the scalar commit transcript
\[
\tau^{(c)}:=\bigl(W_0,\ cr_1,\ W_1^{(c)},\ \dots,\ cr_k,\ W_k^{(c)},\ \gt(c)\bigr),
\qquad W_i^{(c)}(x):=W_i(x)(c).
\]
Its passing sets and passing ratios are denoted $P_i^{(c)}$, $\pi_i^{(c)}$.
\end{definition}

\begin{lemma}[Specialization enlarges passing sets]\label{lem:spec-passing}
For every $c\in\K$ and $0\le i\le k$: $P_i^{\mathrm{sym}}\subseteq P_i^{(c)}$. In
particular $\pi^{\mathrm{sym}}_0\le\pi^{(c)}_0$.
\end{lemma}

\begin{proof}
Descending induction. If $W_k(y)=\gt$ in $\K[T]$ then $W_k^{(c)}(y)=\gt(c)$, so
$P_k^{\mathrm{sym}}\subseteq P_k^{(c)}$. If $y\in Q^{\mathrm{sym}}_{i+1}$, then
$y\in P^{\mathrm{sym}}_{i+1}\subseteq P^{(c)}_{i+1}$ and, by \cref{lem:specialization},
$W^{(c)}_{i+1}(y)=\Fold_{rT}[W_i](y)|_{T=c}=\Fold_{cr_{i+1}}[W_i^{(c)}](y)$, so
$y\in Q^{(c)}_{i+1}$. Taking preimages under $\sigma_{i+1}$ gives
$P^{\mathrm{sym}}_i\subseteq P^{(c)}_i$.
\end{proof}

%=====================================================================
\section{The coupling}
%=====================================================================

\begin{lemma}[Coupling]\label{lem:coupling}
Let $\mathsf P^\ast$ be a symbolic prover. Consider the experiment that samples
$c\leftarrow\K^\times$ uniformly and independently, then runs the symbolic game with
$\mathsf P^\ast$. Define the scalar prover $\mathsf A$ whose coins are $c$ together with
the coins of $\mathsf P^\ast$: it outputs $W_0$; on receiving $\beta_{i+1}$ it sets
$r_{i+1}:=\beta_{i+1}/c$, feeds $r_{i+1}$ to $\mathsf P^\ast$, and outputs
$W_{i+1}^{(c)}$; finally it outputs $\gt(c)$. Then:
\begin{enumerate}[label=\textup{(\roman*)}]
\item the joint law of $(c,\tau)$ in the scalar game with $\mathsf A$ (where $\tau$ is the
symbolic transcript simulated by $\mathsf A$) equals the joint law of $(c,\tau)$ in the
experiment;
\item the scalar commit transcript produced by $\mathsf A$ is $\tau^{(c)}$;
\item in the experiment, $c$ is independent of $\tau$.
\end{enumerate}
\end{lemma}

\begin{proof}
(ii) holds by construction. For (i), condition on $c$, the coins of $\mathsf P^\ast$ and
the history up to $W_i$. In the scalar game $\beta_{i+1}$ is uniform on $\K$ and
independent of this conditioning; since $\beta\mapsto\beta/c$ is a bijection of $\K$,
$r_{i+1}=\beta_{i+1}/c$ is uniform on $\K$ and independent of it as well, which is its law
in the experiment. Induction over $i$ gives (i). (iii) holds by definition of the
experiment.
\end{proof}

%=====================================================================
\section{The main theorem}
%=====================================================================

Recall $\Lambda_\alpha(W_0)$ (\cref{def:block-list}), $\ell_\alpha$, $f_G=\kappa^{-1}(G)$
and the degree cap $d=\sum_i(a_i-1)$. For $G\in\K[X]_{<L}$, $f_G(rT)$ abbreviates
$f_G(r_1T,\dots,r_kT)\in\K[T]$; its degree is at most the total degree of $f_G$, hence at
most $d$.

\begin{theorem}[Symbolic list binding]\label{thm:symbolic-list}
Let $\alpha\in[\rho,1]$ and $\delta:=1-\alpha$. For every symbolic prover,
\[
\Pr\Bigl[\Acc\ \wedge\ \gt\notin\bigl\{f_G(rT):G\in\Lambda_\alpha(W_0)\bigr\}\Bigr]
\ \le\ \eps_{\MCA}(\delta)+\frac{\ell_\alpha\,d}{|\K|-1}+(\alpha^-)^{s},
\]
where $\alpha^-\le\alpha$ is the largest element of $\frac1{n_1}\Z$ smaller than $\alpha$.
\end{theorem}

\begin{proof}
Work in the experiment of \cref{lem:coupling}; adding the independent $c$ does not change
the probability of any event of the symbolic game. Let $\mathsf{Win}$ be the event in
the statement and $\mathsf{Bad}^{(c)}$ the event $\mathsf{Bad}(\delta)$ of
\cref{def:bad-event} for $\tau^{(c)}$. By \cref{lem:coupling}(i),(ii) and
\cref{lem:bad-prob} applied to $\mathsf A$,
$\Pr[\mathsf{Bad}^{(c)}]\le\eps_{\MCA}(\delta)$.

Define, as a function of $\tau$,
\[
Z(\tau):=\bigl\{c'\in\K^\times:\ \exists G\in\Lambda_\alpha(W_0)\ \text{with}\ \gt(c')=f_G(c'r)\bigr\}.
\]
If $\gt\notin\{f_G(rT):G\in\Lambda_\alpha(W_0)\}$, each $\gt-f_G(rT)$ is a nonzero
polynomial of degree $\le d$, so $|Z(\tau)|\le\ell_\alpha d$. As $c$ is uniform on
$\K^\times$ and independent of $\tau$ (\cref{lem:coupling}(iii)),
\[
\Pr\bigl[\mathsf{Win}\wedge c\in Z(\tau)\bigr]\le\frac{\ell_\alpha d}{|\K|-1}.
\]

Now fix $\tau$ and $c$ such that
$\gt\notin\{f_G(rT):G\in\Lambda_\alpha(W_0)\}$, $\neg\mathsf{Bad}^{(c)}$ and
$c\notin Z(\tau)$. Suppose $\pi^{\mathrm{sym}}_0\ge\alpha$. Then
$\pi^{(c)}_0\ge\alpha$ by \cref{lem:spec-passing}, and \cref{lem:master} applied to
$\tau^{(c)}$ yields $V_0\in\Lambda_\alpha(W_0)$ with
$\widehat V_0(cr)=\gt(c)$, i.e.\ $f_{V_0}(cr)=\gt(c)$ by \cref{lem:full-fold}. Here we
used that the initial word of $\tau^{(c)}$ is $W_0$ itself, because $W_0$ does not depend
on $T$. This means $c\in Z(\tau)$, a contradiction. Hence
$\pi^{\mathrm{sym}}_0<\alpha$; since $\pi^{\mathrm{sym}}_0=|Q^{\mathrm{sym}}_1|/n_1\in\frac1{n_1}\Z$,
$\pi^{\mathrm{sym}}_0\le\alpha^-$ and $\Pr[\Acc\mid\tau]\le(\alpha^-)^s$ by
\cref{lem:passing-sym}. Summing the three contributions proves the theorem.
\end{proof}

\begin{definition}[Opened polynomial]\label{def:opened}
For a symbolic commit transcript $\tau$ let $G(\tau)$ be the unique
$G\in\Lambda_\alpha(W_0)$ with $f_G(rT)=\gt$ if exactly one such $G$ exists, and
$G(\tau):=\bot$ otherwise.
\end{definition}

\begin{theorem}[Identification and agreement]\label{thm:identification}
Let $\alpha\in[\rho,1]$, $\delta:=1-\alpha$, and let $\mathsf E$ be the event
\[
\pi^{\mathrm{sym}}_0\ge\alpha\quad\text{and}\quad
\bigl(G(\tau)=\bot\ \ \vee\ \ \pi^{\mathrm{sym}}_0>\agr(W_0,G(\tau))\bigr).
\]
For every symbolic prover,
\[
\Pr[\mathsf E]\le\eps_{\MCA}(\delta)+\Bigl(\ell_\alpha+\binom{\ell_\alpha}2\Bigr)\frac{d}{|\K|-1}.
\]
\end{theorem}

\begin{proof}
Work in the coupled experiment. Let $\mathsf{Inj}$ be the event that
$G\mapsto f_G(rT)$ is injective on $\Lambda_\alpha(W_0)$. For $G\ne G'$ in the list,
$h:=f_G-f_{G'}\ne0$; write $h=\sum_th_t$ with $h_t$ homogeneous of degree $t$ and pick $t$
with $h_t\ne0$. Then $h(rT)=\sum_th_t(r)T^t$, so $h(rT)=0$ forces $h_t(r)=0$, which has
probability at most $t/|\K|\le d/|\K|$ over $r$ uniform in $\K^k$ (Schwartz--Zippel). Since
$\Lambda_\alpha(W_0)$ is fixed before $r$ is sampled,
$\Pr[\neg\mathsf{Inj}]\le\binom{\ell_\alpha}2d/|\K|$.

Assume $\neg\mathsf{Bad}^{(c)}$ and $\pi^{\mathrm{sym}}_0\ge\alpha$. As in the proof of
\cref{thm:symbolic-list}, \cref{lem:spec-passing,lem:master} give
$G_c\in\Lambda_\alpha(W_0)$ with $f_{G_c}(cr)=\gt(c)$ and
$\pi^{\mathrm{sym}}_0\le\pi^{(c)}_0\le\agr(W_0,G_c)$. If moreover $\mathsf{Inj}$ holds
and $G(\tau)=G_c$, then $\mathsf E$ fails. Otherwise either $\neg\mathsf{Inj}$, or
$\mathsf{Inj}$ holds and $c$ is a root of the nonzero polynomial $\gt-f_{G'}(rT)$ for some
$G'\in\Lambda_\alpha(W_0)$ with $f_{G'}(rT)\ne\gt$ (namely $G'=G_c$); the latter set of
$c$ has size at most $\ell_\alpha d$ and is determined by $\tau$, so it has probability at
most $\ell_\alpha d/(|\K|-1)$.
\end{proof}

%=====================================================================
\section{Affine lines and polynomial curves}\label{sec:curves}
%=====================================================================

\begin{definition}[Curve challenges]\label{def:curve-challenges}
Fix $e\ge1$ and public $z=(z_1,\dots,z_k)\in\K^k$. In \cref{prot:sfri} replace the
challenge $r_{i+1}T$ by
\[
\varphi_{i+1}(T):=z_{i+1}+\sum_{t=1}^{e}r_{i+1,t}T^t,\qquad
(r_{i+1,1},\dots,r_{i+1,e})\leftarrow\K^e\ \text{uniform},
\]
raise the degree caps to $d_i^{(e)}:=e\,d_i$, and require $\deg\gt\le ed$. The honest
terminal polynomial is $g(T)=f(\varphi_1(T),\dots,\varphi_k(T))$. The case $e=1$, $z=0$ is
\cref{prot:sfri}; the case $e=1$ certifies $f$ on the affine line
$\{z+cr:c\in\K\}$.
\end{definition}

\begin{proposition}[Curve version]\label{prop:curves}
\Cref{thm:symbolic-list,thm:identification} hold for the protocol of
\cref{def:curve-challenges}, with $f_G(rT)$ replaced by $f_G(\varphi(T))$ and $d$ replaced
by $ed$; so do all results of \cref{ch:corollaries}, and the attacks of
\cref{ch:lower-bounds} apply verbatim with the same replacements. The point $z$ may be
sampled by the verifier after $W_0$, provided it is independent of the $r_{i,t}$.
\end{proposition}

\begin{proof}
We list the steps of \cref{ch:setting} and of this chapter that use the form of the
challenge and check each; conditioning on $z$ does not affect any of them, since $z$ is
independent of the $r_{i,t}$ and $W_0$ is fixed before the $r_{i,t}$ are sampled.
\begin{enumerate}[label=\textup{(\roman*)}]
\item \emph{Honest invariant.} \Cref{lem:fold-poly} holds for every $\beta\in\K[T]$, so
folding with $\varphi_{i+1}(T)$ substitutes $Y_{i+1}=\varphi_{i+1}(T)$ in the proof of
\cref{prop:honest}; the $T$-degree of layer $i$ is at most $e\,d_i$.
\item \emph{Specialization.} The proof of \cref{lem:specialization} gives
$\Fold_{\varphi(T)}[w](y)|_{T=c}=\Fold_{\varphi(c)}[w^{(c)}](y)$.
\item \emph{Coupling.} Fix $c\in\K^\times$. Given $\beta_{i+1}$ uniform, the scalar prover
samples $r_{i+1,2},\dots,r_{i+1,e}$ uniformly and sets
$r_{i+1,1}:=c^{-1}\bigl(\beta_{i+1}-z_{i+1}-\sum_{t\ge2}r_{i+1,t}c^t\bigr)$. For fixed
$(c,r_{i+1,2},\dots,r_{i+1,e})$ this is a bijection between $\beta_{i+1}$ and
$r_{i+1,1}$, so the simulated vector $(r_{i+1,t})_t$ is uniform on $\K^e$ and independent
of the history, and $\varphi_{i+1}(c)=\beta_{i+1}$. \Cref{lem:coupling} follows as before.
\item \emph{Root count.} $\gt-f_G(\varphi(T))$ has degree at most $ed$, so
$|Z(\tau)|\le\ell_\alpha ed$.
\item \emph{Injectivity.} For $0\ne h\in\Pa(\K)$, if $h(\varphi(T))=0$ in $\K[T]$ then
$h(\varphi(1))=0$, where $\varphi(1)=z+\sum_tr_{\cdot,t}$ is uniform on $\K^k$; this has
probability at most $\deg h/|\K|\le d/|\K|$ by Schwartz--Zippel. This replaces the
homogeneous-component argument in the proofs of \cref{thm:identification} and
\cref{lem:perturbation}(iv).
\end{enumerate}
No other step depends on the challenge.
\end{proof}

%=====================================================================
\section{Remarks on the protocol}
%=====================================================================

\begin{remark}[Which degree checks matter]\label{rem:degree-checks}
The analysis uses exactly two syntactic checks.
\begin{enumerate}[label=\textup{(\roman*)}]
\item \emph{$W_0$ is $T$-free.} This makes the initial word of $\tau^{(c)}$, hence the
list $\Lambda_\alpha(W_0)$ and the reference $F_\star(W_0)$, independent of $c$. Without
it the root count in $c$ is not available.
\item \emph{$\deg\gt\le d$.} This bounds the number of roots of $\gt-f_G(rT)$. A cap
$d'\ge d$ would replace $d$ by $d'$ in all bounds.
\end{enumerate}
The caps $d_i$ for $1\le i<k$ only fix the encoding length of the oracles.
\end{remark}

\begin{remark}[Checking at a single random point]\label{rem:single-point}
Consider the variant of \cref{prot:sfri} in which, after $\gt$ is sent, $\mathsf V$
samples $c'\leftarrow\K^\times$ and performs every check Q2--Q3 only at $T=c'$. The proof
of \cref{thm:symbolic-list} applies with $c:=c'$: acceptance now has conditional
probability $(\pi^{(c')}_0)^s$, the point $c'$ is independent of $\tau$, and outside
$\mathsf{Bad}^{(c')}\cup\{c'\in Z(\tau)\}$ the master path lemma forces
$\pi^{(c')}_0<\alpha$. Hence the bound of \cref{thm:symbolic-list} holds verbatim for
this variant. The opened data per query are unchanged; only the verifier's arithmetic
changes.
\end{remark}

\begin{remark}[Reuse of evaluations]\label{rem:reuse}
Outside the event of \cref{thm:symbolic-list}, $\gt=f_G(rT)$ as polynomials for some
$G\in\Lambda_\alpha(W_0)$. Then $\gt(c'')=f_G(c''r_1,\dots,c''r_k)$ holds for all
$c''\in\K$ simultaneously, with no additional error per evaluation and no union bound
over the number of evaluations. The certified set is the line
$\{c''r:c''\in\K\}\subseteq\K^k$ through the origin in the direction $r$.
\end{remark}
